\documentclass[11pt]{article}

\usepackage[utf8]{inputenc}
\usepackage{amsmath}
\usepackage{graphicx}
\usepackage{hyperref}
\usepackage{cite}

\title{A Minimal Calibration Paper Skeleton}
\author{Workshop Attendee \\ \texttt{you@example.org}}
\date{\today}

\begin{document}
\maketitle

\begin{abstract}
This is a placeholder abstract for the LaTeX exercise. It describes the
synthetic instrument-calibration workflow used elsewhere in this workshop
repository, with just enough scaffolding to demonstrate citations,
cross-references, equations, and figures.
\end{abstract}

\section{Introduction}
\label{sec:intro}

Modern atmospheric instruments require careful per-channel calibration to
remove systematic offsets and scale factors~\cite{anderson2024calibration}.
Recent work has emphasized robust statistical methods for combining
calibration constants across multiple deployments~\cite{nguyen2023robust}.

In this paper we describe a simple linear calibration applied to a
three-channel synthetic dataset, demonstrating the workflow used by the
VSCode workshop exercises.

\section{Methods}
\label{sec:methods}

For each observation, the calibrated value $y$ is obtained from the raw
value $x$ via the per-channel linear transform
\begin{equation}
    y_c = s_c \, x_c + o_c,
    \label{eq:calibration}
\end{equation}
where $s_c$ and $o_c$ are the slope and offset for channel $c$.

Summary statistics over the calibrated dataset $\{y_i\}_{i=1}^{N}$ are
computed as the mean, minimum, and maximum:
\begin{align}
    \bar{y} &= \frac{1}{N} \sum_{i=1}^{N} y_i, \\
    y_{\min} &= \min_i y_i, \\
    y_{\max} &= \max_i y_i.
\end{align}

\section{Results}
\label{sec:results}

Figure~\ref{fig:calibration} shows the raw and calibrated distributions
side-by-side. Cross-reference: see Section~\ref{sec:methods} for the
underlying formula (Equation~\ref{eq:calibration}).

\begin{figure}[h]
    \centering
    % To make this render, drop a real PNG in this folder and point at it.
    % e.g. ../../self-study/02-jupyter/calibration_plot.png (after running the notebook)
    \fbox{\rule{0pt}{6em}\rule{12em}{0pt}\textit{[figure placeholder]}}
    \caption{Raw versus calibrated distributions for three channels.}
    \label{fig:calibration}
\end{figure}

\section{Discussion}
\label{sec:discussion}

The linear calibration of Section~\ref{sec:methods} is sufficient for the
synthetic case here, but real instruments often require higher-order or
piecewise corrections.

\bibliographystyle{plain}
\bibliography{refs}

\end{document}
